\section{Fourth moments of singular bands}
\label{rec:sec-bands}
\label{sec:band-grid}

We now retain the carrier rank of each local singular band directly
in a fourth moment. The overlap estimate applies to local factors
with a specified Fourier type, rank and operator norm. Raising the
child factors to a suitable power cancels the rank cost, and
interpolation returns to the original sweep. The resulting indexed
power bound is equivalent, after changing an absolute exponent, to
Theorem~\ref{rec:perspective-singular}; this proof does not carry its
level--dimension weight.

\subsection{The sparse estimate from coordinate chunks}

The sparse argument uses a complete sweep followed by a partially
coupled reversed sweep. Cancellation forces every tracked label to
meet another during the second sweep. A fixed number of coordinate
chunks converts this coverage event into an occupation bound supplied
by the first sweep.

Here the sparse input works up to any fixed power less than one.
For each fixed $0<\delta<1$, there are $c_\delta>0$ and
$n_\delta<\infty$ such that, for every dyadic $n\ge n_\delta$
and every $\lambda\vdash n$,
\begin{equation}\label{rec:chunk-sparse}
 0\le k=n-\lambda_1\le n^{1-\delta}
 \quad\Longrightarrow\quad
 \|T_n|_{V_\lambda}\|_{\op}\le e^{-c_\delta k\log n}.
\end{equation}
We will use this with $\delta=1/200$.
\begin{proof}
For $k=0$ the type is trivial and the assertion is $1\le1$.
We specify the tuple kernels and their normalization first.
Let $\mathcal X_k$ be the ordered injective $k$-tuples of sites.
Write $P(x,y)$ for the forward sweep transition probability on
this space, with input index $x$ and output index $y$.
As an operator on tuple functions this is the adjoint of the
homomorphic sweep action, and has the same singular values
on every irreducible block.  For $I\subseteq[k]$, define
\begin{equation}\label{rec:sparse-partial-kernel}
 Q_I(x,y)=n^{-(k-|I|)}
              \Pr\{g(x_I)=y_I\},\qquad x,y\in\mathcal X_k.
\end{equation}
Before restriction to injective outputs, this kernel moves the
$I$-labels with one common sweep and moves every other label
by its own independent fair bit at each coordinate.
The latter endpoints are independent uniform sites.
Thus $Q_I$ has row sums at most $1$.
Its transpose has the same description using the reversed
sweep, and has row sums at most $1$ as well.
All sums here are counting sums of transition probabilities.
Rescaling the tuple inner product to uniform probability does
not change the operator or its norm.

The block $\lambda=(n-k,\beta)$ occurs on $\mathcal X_k$ and
does not occur on any proper subset of its slots.
For $I\ne[k]$, both the range of $Q_I$ and the range of $Q_I^*$
consist of functions of the $I$-coordinates.  These ranges
have no $\lambda$-part by branching.  Consequently
\begin{equation}\label{rec:sparse-centered-kernel}
 Z=\sum_{I\subseteq[k]}(-1)^{k-|I|}Q_I
 \quad\hbox{satisfies}\quad
 Z|_\lambda=P|_\lambda,\qquad Z^*|_\lambda=P^*|_\lambda.
\end{equation}
Here and below a restriction to $\lambda$ may include all its
copies in the tuple module.

Given endpoints $x,y$, the path of each label through a sweep
is unique: its updated coordinates have their values in $y$
and its remaining coordinates have their values in $x$.
Join two labels when these paths use the same switch at some
time, including an inconsistent use of the same switch input.
Let $\Gamma(x,y)$ mean that this graph has no isolated vertex.
If a label is isolated, its prescribed switches are independent
of those prescribed by all the other labels.  Adding this label
to $I$ therefore multiplies its joint endpoint probability by
$1/n$, exactly canceled by the change in the prefactor in
\eqref{rec:sparse-partial-kernel}.  Pairing terms with and
without that label proves
\begin{equation}\label{rec:sparse-cover-majorant}
 |Z(x,y)|\le\mathbf1_{\Gamma(x,y)}
                         \sum_{I\subseteq[k]}Q_I(x,y).
\end{equation}
For $k=1$ the two terms are identical, so $Z=0$ and the
assertion follows.  Henceforth take $k\ge2$.

On the $\lambda$-part, $PZ^*=PP^*$ is positive semidefinite.
An absolute row-sum bound for $PZ^*$ therefore bounds
$\|P|_\lambda\|_{\op}^2$: apply the operator to an eigenvector
of $PP^*|_\lambda$ and choose a coordinate of maximum
absolute value.  By \eqref{rec:sparse-cover-majorant}, it
suffices to bound, uniformly in an initial injective tuple,
the following experiment for each $I$.  Run a full sweep,
obtaining an injective intermediate tuple $z$; then run a
partially coupled reversed sweep, jointly on $I$ and
independently on the other labels, and ask for $\Gamma$ in
this second sweep.  Dropping injectivity of the final tuple
only increases the nonnegative bound.  This description
follows exactly by multiplying $P(x,z)Q_I(y,z)$ and summing
over $z,y$; no replacement of the two-sweep law is made.

The occupied set of $z$ satisfies
\begin{equation}\label{rec:chunk-intermediate-occupation}
 \Pr\{A\subseteq\{z_1,\ldots,z_k\}\}
 \le(k/n)^{|A|}
\end{equation}
for every fixed set of sites $A$.  This is the case $r=k$,
$s=0$ of Lemma~\ref{grid:occupation}, applied to singleton
test sets.  Its hypotheses hold because this first sweep
moves all $k$ labels jointly from distinct sites.

Put $L=\lfloor\delta d/2\rfloor$ and partition the second
sweep's chronological coordinate list into $J=\lceil d/L\rceil$
nonempty consecutive chunks, each of length at most $L$.
For sufficiently large $d$, $L\ge1$ and
$J\le J_\delta:=\lceil4/\delta\rceil+1$.
A chunk of length $\ell$ partitions sites into bins of size
$s=2^\ell\le n^{\delta/2}$, according to all the coordinates
outside that chunk.  Each path stays in its bin throughout
the chunk.  Every contact therefore lies in a bin containing
at least two labels at the start of that chunk.

On $\Gamma$, every label is in such a bin in at least one
chunk.  Thus some chunk has at least $k/J$ labels in bins
of occupancy at least two.  Define
\[
 r=\max\{1,\lfloor k/(2J)\rfloor\}.
\]
If $k\ge2J$, then $2r\le k/J$.  If there are at least $r$
multiply occupied bins, take any $r$ of them.  Otherwise take
all of them and pad to $r$ bins: their total occupancy is
still at least $k/J\ge2r$.  If $k<2J$, coverage supplies
at least one contact, so some chunk has one bin with two
labels and the same conclusion holds with $r=1$.
There are enough bins for padding, since
$n/s\ge n^{1-\delta/2}>k\ge r$ for large $n$.
In all cases
\begin{equation}\label{rec:chunk-selected-bins}
 r\ge k/(4J),\qquad
 \Gamma\ \Longrightarrow\
 \text{some chunk has $r$ bins containing at least $2r$ labels}.
\end{equation}

Fix a chunk and a specified union $A$ of $r$ of its bins.
Then $|A|=rs$.  Presample the partially coupled motion
from the beginning of the second sweep to this chunk.
Use one common random permutation $\phi$ for the $I$-labels.
Independently, for every possible starting site $v$, sample
one independent walk with endpoint $\xi_v$ after that prefix.
If $z_i=v$ and $i\notin I$, use this walk for label $i$.
The $z_i$ are distinct, so this site-indexed construction
gives exactly independent walks for those labels.
All these random maps are independent of the first sweep.

Put $C=\phi^{-1}(A)$ and $D=\{v:\xi_v\in A\}$.
The number of labels in $A$ at the start of the chunk is
at most the number of occupied intermediate sites in $C\cup D$.
We have $|C|=rs$, while $W=|D|$ is a sum of independent
Bernoulli variables with mean $rs$.  Indeed the transition
matrix of any coordinate prefix is doubly stochastic, so
$\sum_v\Pr(\xi_v\in A)=|A|$.
Consequently
$\E\binom Wj\le(rs)^j/j!$.
Conditioning on the presampled maps and using
\eqref{rec:chunk-intermediate-occupation}, then averaging,
gives
\[
 \begin{split}
 \Pr\{\text{at least $2r$ labels in }A\}
 &\le(k/n)^{2r}\E\binom{rs+W}{2r}\\
 &\le(k/n)^{2r}
       \sum_{j=0}^{2r}\frac{(rs)^{2r-j}}{(2r-j)!}
                            \frac{(rs)^j}{j!}\\
 &=(k/n)^{2r}\frac{(2rs)^{2r}}{(2r)!}.
 \end{split}
\]
This is an unconditional preimage calculation.  It does not
condition the intermediate tuple on previous contacts.
Summing over choices of the $r$ bins yields
\[
 \binom{n/s}{r}(k/n)^{2r}\frac{(2rs)^{2r}}{(2r)!}
 \le\left(\frac{e^3k^2s}{nr}\right)^r
 \le(CJks/n)^r.
\]
Now $ks/n\le n^{-\delta/2}$.  For all sufficiently large
$n$, depending only on $\delta$, this last bound is at most
$n^{-\delta r/4}\le n^{-\delta k/(16J_\delta)}$.
Union over the $J$ chunks and the $2^k$ choices of $I$
therefore gives
\[
 \sup_x\sum_y|(PZ^*)(x,y)|
 \le2^kJ_\delta n^{-\delta k/(16J_\delta)}
 \le n^{-\delta k/(32J_\delta)}.
\]
The positive-block eigenvalue argument above and a square
root prove \eqref{rec:chunk-sparse}, for example with
$c_\delta=\delta/(64J_\delta)$ after increasing the
starting size.  The separate $k=1$ argument supplies the
smallest level without any padding or rounding exception.
\end{proof}

\subsection{A fourth moment with local ranks}

\begin{lemma}\label{rec:band-grid-entropy}
Put $m=\min(a,b)$ and assume $a,b\le2m$, with $m$ sufficiently
large. Write $K=(S_b)^a$ and $L=(S_a)^b$ for the row and column
groups. Let $Y=Y_LY_K$, where $Y_K$ and $Y_L$ are tensor products
of operators on the row and column permutation groups. Assume each
local factor is supported in one irreducible Fourier matrix of
dimension $D_i$, has rank at most $r_i$, and has operator norm at
most $b_i$. If a local factor is zero, then $Y=0$. Otherwise
$r_i>0$ for every $i$; put $x=\sum_i\log(D_ir_i)$. In this case,
for every
$\lambda\vdash n=ab$,
\[
 D_\lambda\Tr|Y_\lambda|^4
 \le \exp\{Cn m^{-1/16}+Cx/\log m\}
                         \prod_i D_ir_i b_i^4.
\]
\end{lemma}
\begin{proof}
The combinatorial assertion needed after expanding matrix coefficients
is a bound on compatibility between biased row and column permutations:
\[
 \frac{|S_n|}{|K||L|}\Pr(\text{compatible})
 \le\exp\{Cn m^{-1/16}+C(x_K+x_L)/\log m\}.
\]
Here a row law has density at most $e^{x_i}$ and a column law $q_c$
has density at most $e^{x'_c}$ relative to the appropriate uniform
permutation law; set $x_K=\sum_i x_i$ and $x_L=\sum_cx'_c$.
Fix the row permutations $\sigma_i$.  Form a bipartite multigraph
whose left vertices are output columns $c$ and whose right vertices
are original columns $j$, with an edge of index $i$ when
$\sigma_i(j)=c$.  Its degree is $a$ and its edge multiplicities are
$m_{cj}$.  The column permutation $\tau_c$ assigns color $\tau_c(i)$
to that edge.  Compatibility is exactly the requirement that these
colors are proper also at the right vertices.  Put
$H=\sum_{c,j:\,m_{cj}\ge m^{1/2}}m_{cj}$.

Let $p$ be the probability of proper coloring under $\otimes_cq_c$.
For $p>0$, let $\nu$ be that law conditioned on proper coloring,
with column marginals $\nu_c$.  Write $\mathcal H$ for Shannon
entropy, and define
\begin{equation}\label{rec:band-conditioned-KL}
 \begin{aligned}
 h_c&=\log(a!)-\mathcal H(\nu_c)=\KL(\nu_c\Vert U_c),\\
 I&=\sum_c\mathcal H(\nu_c)-\mathcal H(\nu),
 &D&=\sum_c\KL(\nu_c\Vert q_c).
 \end{aligned}
\end{equation}
Here $U_c$ is uniform on the $a!$ column permutations.
The same chain rule as in \eqref{rec:biased-success-KL} gives
$-\log p=I+D$ and $\sum_ch_c\le D+x_L$.

Reveal colors in order, and for each color expose its edges in an
independent random priority order of the left vertices $c$.
Let $p_c$ be the prediction of its edge under the marginal $\nu_c$
given its own earlier colors, and let $P_j$ be the resulting target
distribution.  The chain-rule sum of the true predictions' expected
KL divergences from these marginal predictions is $I$.
At a successful coloring, only targets unused by earlier vertices
in this color are allowed.  Conditional on the coloring and current
priority, each target other than the actual target $j^*$ belongs to
one other vertex in the color matching.  It remains available with
probability $t$, where $t$ is uniform on $[0,1]$.  Jensen therefore
gives the integrated lower bound
\[
 \int_0^1-\log(t+(1-t)P_{j^*})\,dt
 =1-\ell(P_{j^*}),\qquad
 \ell(z)=\frac{-z\log z}{1-z}\le\min(1,C\sqrt z),
\]
with continuous endpoint values.

The last $\lceil m^{3/4}\rceil$ colors cost at most $bm^{3/4}$
in lost unit contributions.  Targets with multiplicity at least
$m^{1/2}$ cost at most $H$.  At any other exposure, with at least
$v\ge m^{3/4}$ colors remaining, a light target has mass at most
$m^{-1/4}$ under the image of the uniform prediction on the $v$
remaining edges.  If $P_j<m^{-1/8}$, its loss is at most
$Cm^{-1/16}$.  For the other light targets let $z$ be their total
$P$-mass.  Projecting the prediction divergence
$e=\log v-\mathcal H(p_c)$ onto the target distribution and using
log-sum on the complement gives
$e\ge z(\log m)/8-z$.  Hence $z\le Ce/\log m$ for large $m$.
The expected sum of these prediction divergences at vertex $c$ is
exactly $h_c$ by the marginal entropy chain rule.
Summing all losses proves
\begin{equation}\label{rec:band-total-correlation}
 I\ge n-Cn m^{-1/16}-H-a_m\sum_ch_c,\qquad
 a_m=C/\log m.
\end{equation}
Adding $D$ and using \eqref{rec:band-conditioned-KL} now gives,
with the same positive remainder as in
\eqref{rec:biased-entropy-absorption},
\begin{equation}\label{rec:band-entropy-absorption}
 -\log p\ge n-Cn m^{-1/16}-H-a_m x_L+(1-a_m)D.
\end{equation}
For large $m$, discard the last term.  Thus, conditionally on the
rows, $p\le\exp(-n+Cn m^{-1/16}+H+Cx_L/\log m)$.
When $p=0$ this bound is automatic.

Under uniform independent row permutations, a witness count gives,
for an absolute $\eta>0$,
\[
 \E\exp(\eta\log m\,H)\le2.
\]
Indeed, if $H=h>0$, its heavy cells number at most $h/m^{1/2}$.
Choosing those cells and their row incidences, and then specifying
the row-permutation images, gives probability at most
$h b^{2h/m^{1/2}}(e^2a/(bm^{1/2}))^h\le e^{-ch\log m}$.
Here a consistent specification of $u$ images in one row has
probability $1/(b)_u\le(e/b)^u$; inconsistent specifications have
probability zero.  Summing this tail proves the displayed moment.
The biased row density has mean one and is bounded by $e^{x_K}$.
H\"older with exponent $\eta\log m$ therefore gives
$\E_{\rm biased}e^H\le(2e^{x_K})^{1/(\eta\log m)}$.
Average the conditional success bound and use the factorial
normalization $\log(|S_n|/(|K||L|))\le n+C\log n$.
The exact $-n$ term cancels, leaving the asserted compatibility
bound.

To pass to arbitrary local factors, put $X_L=Y_L^*Y_L$ and
$X_K=Y_KY_K^*$. Cyclicity identifies $\Tr|Y|^4$ with
$\Tr(X_LX_KX_LX_K)$. A local positive factor has regular rank
at most $D_ir_i$ and regular squared Hilbert--Schmidt norm at most
$D_ir_ib_i^4$. Lemma~\ref{lem:coefficient-compatibility} therefore
bounds the regular four-factor trace by the compatibility probability
just estimated, multiplied by $\prod_iD_ir_ib_i^4$. The lemma's
squared-coefficient densities have caps at most $D_ir_i$, so the
parameter in that probability estimate is $x=\sum_i\log(D_ir_i)$.
Each irreducible contribution is nonnegative and has multiplicity
$D_\lambda$, proving the displayed bound.
\end{proof}

\subsection{Closing the singular-band recursion}

\begin{theorem}\label{rec:all-rank-power}
There is an absolute $\alpha>0$ such that
$s_t(T_n|_{V_\lambda})\le(D_\lambda t)^{-\alpha}$ for every
dyadic $n\ge2$, every $\lambda\vdash n$, and every
$1\le t\le D_\lambda$.
\end{theorem}
\begin{proof}
Lemma~\ref{lem:positive-power-comparison} gives, for square
matrices $A_1,A_2$ and real $P\ge q\ge1$,
\begin{equation}\label{rec:band-power-interpolation}
 \|A_2A_1\|_{S^P}
 \le\big\||A_2|^{P/q}|A_1^*|^{P/q}\big\|_{S^q}^{q/P}.
\end{equation}
This applies to noninvertible factors and uses unnormalized norms.

We now prove the recursive estimate.  Suppose that the two smaller
sweeps satisfy the theorem with a common exponent
$0<\bar\alpha\le1/8$.  A smaller exponent weakens the assertion,
so later we will take the minimum of the two inherited exponents.
Write $m=\min(a,b)$, $\ell=\log m$, and
\[
 p=\frac{1+K_0/\ell}{4\bar\alpha},
\]
where the absolute constant $K_0$ will be fixed below.
In particular $p\ge2$.
Recall the row-first factorization $T_n=YX$.
The positive matrices $|Y|$ and $|X^*|$ are tensor products of
the corresponding local absolute values on the column and row
groups.  Each local absolute value has exactly the singular
values of a child sweep.

Here is a finite band decomposition whose bounds do not depend
on how small $\bar\alpha$ is.  On a local type
$\rho\vdash q$, with $q=a$ or $b$, put $D=D_\rho$ and list
the positive eigenvalues of its absolute value $S_\rho$ in
nonincreasing order:
$\sigma_1\ge\cdots\ge\sigma_R>0$, where $R\le D$.
Fix an orthonormal eigenbasis, including an arbitrary choice
within a repeated eigenspace.  For each nonempty index block
\[
 I_v=\{2^v,\ldots,\min(2^{v+1}-1,R)\},\qquad v\ge0,
\]
let $E_{\rho,v}$ be the projection onto those eigenvectors and
let $r_{\rho,v}=|I_v|$.  Then $r_{\rho,v}\le2^v$, and the
inductive estimate gives the exact bound
\begin{equation}\label{rec:band-local-rank-bound}
 \|S_\rho^pE_{\rho,v}\|_{\op}
 \le(D_\rho\,2^v)^{-p\bar\alpha}
 \le(D_\rho r_{\rho,v})^{-p\bar\alpha}.
\end{equation}
There is no contribution from a zero eigenspace.  This also
handles a local trivial type: its only eigenvalue is $1$ and
$D_\rho r_{\rho,0}=1$.  A local type killed by the child sweep
has no bands.

Fourier inversion realizes each $S_\rho^pE_{\rho,v}$ as a local
group-algebra operator supported in that one Fourier matrix,
zero in every other type.  Summing these operators over types
and bands gives the local positive power.  Tensoring over lines
and distributing the two products therefore expresses
$|Y|^p|X^*|^p$ as a finite sum of operators $Z_\pi$.
A profile $\pi$ chooses a type and one nonempty band on each
of the $a+b$ lines.  No restriction multiplicity has to be
counted separately: each such local Fourier operator already
acts on all its multiplicity copies when extended to $S_n$.
Profiles whose extensions vanish on $V_\lambda$ may be
discarded, and keeping them only enlarges the count below.

Let $\mathfrak p(q)$ denote the number of partitions of $q$.
The partition generating product, evaluated at $e^{-q^{-1/2}}$,
gives $\mathfrak p(q)\le e^{C\sqrt q}$.  Also
$D_\rho\le q!$, so the number of nonempty index bands in any
type is at most $1+\log_2D_\rho\le Cq\log(q+1)$.
Consequently the number $N_q$ of type-band choices on a line
of size $q$ satisfies
\[
 N_q\le e^{C\sqrt q}\,Cq\log(q+1).
\]
Since $m\le a,b\le2m$, the number $\mathcal N$ of profiles obeys
\begin{equation}\label{rec:band-profile-count}
 \log\mathcal N
 \le a\log N_b+b\log N_a
 \le C m^{3/2}\log m.
\end{equation}
The elementary partition bound used here can also be seen by
expanding the logarithm of the generating product:
$\sum_{j,r\ge1}e^{-jr/\sqrt q}/r
\le\sqrt q\sum_{r\ge1}r^{-2}$.
In particular every constant in \eqref{rec:band-profile-count}
is independent of $\bar\alpha$.

For a profile put $x_\pi=\sum_i\log(D_ir_i)$.
Lemma~\ref{rec:band-grid-entropy} and
\eqref{rec:band-local-rank-bound} give
\[
 \begin{split}
 D_\lambda\|Z_\pi|_{V_\lambda}\|_{S^4}^4
 &\le
 \exp\{C_0n m^{-1/16}
       +(1+C_0/\ell-4p\bar\alpha)x_\pi\}\\
 &=
 \exp\{C_0n m^{-1/16}
       +(C_0-K_0)x_\pi/\ell\}.
 \end{split}
\]
Fix $K_0\ge C_0+1$.  Because $x_\pi\ge0$, the last
quantity is at most $\exp(C_0n m^{-1/16})$.
Summation uses the triangle inequality for $S^4$, rather
than an unjustified addition of fourth powers:
\[
 \begin{split}
 \big\|(|Y|^p|X^*|^p)|_{V_\lambda}\big\|_{S^4}
 &\le\sum_\pi\|Z_\pi|_{V_\lambda}\|_{S^4}\\
 &\le\mathcal N D_\lambda^{-1/4}
                  \exp(C_0n m^{-1/16}/4).
 \end{split}
\]
For all sufficiently large $m$, the counting error in
\eqref{rec:band-profile-count} is smaller than a constant
multiple of $n m^{-1/16}$.  Thus, with a fixed absolute
$C_1$ and $E_n=C_1n m^{-1/16}$,
\begin{equation}\label{rec:band-summed-fourth}
 D_\lambda
 \big\|(|Y|^p|X^*|^p)|_{V_\lambda}\big\|_{S^4}^4
 \le e^{E_n}.
\end{equation}
Indeed $4\log\mathcal N=O(m^{3/2}\log m)$, whereas
$n m^{-1/16}\ge m^{31/16}$.  None of the thresholds or
constants in this step depends on $p$ or $\bar\alpha$.

Apply \eqref{rec:band-power-interpolation} on $V_\lambda$,
with $P=4p$ and $q=4$.  Equation
\eqref{rec:band-summed-fourth}, with
$T_\lambda=T_n|_{V_\lambda}$, yields
\[
 \Tr|T_\lambda|^{4p}\le e^{E_n}/D_\lambda.
\]
For each $1\le t\le D_\lambda$, decreasing order of the
singular values gives
$t\,s_t(T_n|_{V_\lambda})^{4p}
\le\Tr|T_\lambda|^{4p}$.
We have proved
\begin{equation}\label{rec:band-interpolation-loss}
 s_t(T_n|_{V_\lambda})
 \le
 \exp\left\{-\frac{\bar\alpha}{1+K_0/\log m}
                  \big(\log(D_\lambda t)-E_n\big)\right\}.
\end{equation}
If the singular value is zero the inequality is automatic.
This argument derives the factor $t$ from a Schatten moment;
it does not require the sweep to be normal.

We give the dimension estimate needed to absorb $E_n$.
For every shape with at most $n/2$ rows,
\begin{equation}\label{rec:band-dimension-lower}
 \log D_\lambda\ge c_1 k,\qquad k=n-\lambda_1,
\end{equation}
for an absolute $c_1>0$ and all sufficiently large $n$.
To check it, put $r=\max(\lambda_1,\lambda'_1)$ and transpose
temporarily if necessary.  If $r\le n/8$, all hooks have
length at most $2r$, so the hook formula gives
$D_\lambda\ge n!/(2r)^n\ge e^{c n}$.
If $n/8<r\le3n/4$, retain a first row of length $r$ and
$j=\lfloor r/4\rfloor$ boxes below it.  Fill its first $j$
top-row boxes first, and then freely interleave the remaining
top row with a fixed standard order on the lower diagram.
Its width is at most $j$, so these are standard tableaux,
and there are $\binom rj\ge e^{c r}$ of them.
They extend to the full diagram.  Finally, $r>3n/4$ must
mean $r=\lambda_1$, because $\lambda'_1\le n/2$.
Now $k<n/4$ and the same construction with $j=k$ gives
$D_\lambda\ge\binom rk\ge(r/k)^k\ge3^k$.
These three cases prove \eqref{rec:band-dimension-lower};
the case $k=0$ requires no lower bound.

For the dense range $k>n^{1-1/200}$ and a shape not killed
by the matching average, let $L=\log(D_\lambda t)$.
Then $L\ge c_1n^{1-1/200}$.  Since $m^2\le n\le2m^2$,
\[
 \frac{E_n}{L}
 \le C n^{\,1/200-1/32}
 =C n^{-21/800}\le\frac1d
\]
for all sufficiently large $d$, with a threshold independent
of $\bar\alpha$.  For $d\ge3$ also
$K_0/\log m\le c_2/d$, where $c_2=3K_0/\log2$.
It follows that
\[
 \frac{1-E_n/L}{1+K_0/\log m}
 \ge\frac{1-1/d}{1+c_2/d}
 \ge1-\frac{C_*}{d},\qquad C_*=1+c_2.
\]
Thus \eqref{rec:band-interpolation-loss} proves the
inductive claim in this range with exponent
$\bar\alpha(1-C_*/d)$.

For $1\le k\le n^{1-1/200}$, write $c_{\mathrm s}>0$
for the fixed constant in \eqref{rec:chunk-sparse}.
The elementary bound $D_\lambda\le n^k$ and $t\le D_\lambda$
give $L\le2k\log n$.  Consequently
\[
 s_t(T_n|_{V_\lambda})
 \le e^{-c_{\mathrm s}k\log n}
 \le (D_\lambda t)^{-c_{\mathrm s}/2}.
\]
This handles the sparse range whenever the exponent being
proved is at most $c_{\mathrm s}/2$.
The trivial representation has $D_\lambda=t=1$ and sweep
equal to $1$, so the theorem is equality there.
The sign representation is zero.  Every other shape killed
by a matching layer, and every further null block, satisfies
the assertion with every positive exponent.

It remains to choose the base and verify a uniform positive
exponent.  Choose an absolute integer $d_0>2C_*$ so large
that every large-size estimate just used holds for $d>d_0$,
including the sparse input, the line-group entropy estimate,
the profile absorption, and $E_n/L\le1/d$.
This choice has been made without specifying an inherited
exponent.  Lemma~\ref{lem:finitegap} then supplies a common positive base
exponent $\alpha_0\le\min(1/8,c_{\mathrm s}/2)$ at the finitely
many sizes $1\le d\le d_0$, exactly as in the finite-base step of
Theorem~\ref{rec:perspective-singular}. Indeed each nonzero,
nonconstant block has norm below one and only finitely many ranks;
zero blocks impose no restriction. The trivial and sign blocks
were handled above.

Set $\alpha_d=\alpha_0$ for $d\le d_0$, and, for $d>d_0$,
define
\[
 \alpha_d=
 \left(1-\frac{C_*}{d}\right)
 \min\{\alpha_{\lfloor d/2\rfloor},
                   \alpha_{\lceil d/2\rceil}\}.
\]
The preceding argument proves the theorem with $\alpha_d$
by induction.  All these numbers are at most $\alpha_0$,
so both the condition $\bar\alpha\le1/8$ in interpolation
and the sparse comparison remain valid.
The minimum selects a branch of rounded halvings, and all
multipliers $1-C_*/d$ are at least $1/2$. By
Lemma~\ref{lem:dyadic-exponents} with $\gamma=1$, their products
are bounded away from zero. Thus $\alpha:=\inf_d\alpha_d>0$
is a common exponent for every type and rank.
For example, a fixed integer $R$ with $2\alpha R\ge3$
then gives
$D_\lambda\|T_n(\lambda)^R\|_{\HS}^2
\le D_\lambda^{\,2-2\alpha R}\le D_\lambda^{-1}$,
the inverse-dimension estimate used in the common mixing
conversion.
\end{proof}
